
\appendix
\crefalias{section}{appendix}
\crefalias{subsection}{appendix}

\renewcommand{\thesection}{\Alph{section}}
\renewcommand{\thesubsection}{\thesection.\arabic{subsection}}
\renewcommand{\theHsection}{\Alph{section}}
\renewcommand{\theHsubsection}{\theHsection.\arabic{subsection}}

\section{Relative nilpotency for groups}
\label{appendix: relative nilpotency for groups}

In this appendix, we collect the relative theory of nilpotency for a normal subgroup $N$ of a group $G$, or equivalently for the quotient epimorphism $G\twoheadrightarrow G/N$.
Although the underlying notions and basic characterizations are classical, we include the proofs in order to fix our conventions and to record statements needed in this paper.
For general background on nilpotent groups, central series, and chief series, see \cite{book_robinson_1982_a_course_in_the_theory_of_groups}.


\begin{HisNote}
The material in this section consists mostly of standard facts that appear in various forms throughout the literature. 
To the best of the authors' knowledge, the earliest explicit formulation we have found of nilpotency for a pair of groups in terms of a relative central series is due to Hilton \cite{hilton_1976_nilpotency_in_group_theory_and_topology,hilton_1982_nilpotency_in_group_theory_and_topology}.

Earlier relative central-series constructions occur in Baer's study of upper central series and hypercentral normal subgroups (\cite{baer_1953_the_hypercenter_of_a_group}). In a different, action-theoretic setting, Kaloujnine~\cite{kaloujnine_1953_uber_gewisse_beziehungen_zwischen_einer_gruppe_und_ihren_automorphismen} introduced $A$-central series and $A$-nilpotency for groups equipped with a group $A$ of automorphisms; in this line of work, Mohamed~\cite{mohamed_1963_on_series_of_subgroups_related_to_groups_of_automorphisms} treated both lower and upper $A$-central series, and Bechtell~\cite{bechtell_1966_frattini_subgroups_and_phi_central_groups} applied Kaloujnine's construction to normal subgroups.
If $N\trianglelefteq G$ and $A$ is the image of the conjugation homomorphism $G\to\Aut(N)$, these operator-central series and the notion of nilpotency specialize to the relative ones considered below.

A further generalization replaces the commutator word by a system of group words; the resulting generalized central series are called \emph{marginal series}. Hulse--Lennox~\cite{hulse_lennox_1976_marginal_series_in_groups} and Fung~\cite{fung_1977_some_theorems_of_hall_type} developed this theory in the 1970s.
For pairs of groups, corresponding relative marginal-series and $\mathscr{V}_G$-nilpotency theories were developed in \cite{moghaddam_salemkar_2002_some_properties_of_ultra_hall_and_schur_pairs,rismanchian_araskhan_2013_some_properties_on_the_baerinvariant_of_a_pair_of_groups_and_mathscrvgmarginal_series}, where $\mathscr{V}$ is a variety of groups. The ordinary relative nilpotency considered here is recovered when $\mathscr{V}=\Ab$. See also \cite{rismanchian_2021_group_extensions_and_marginal_series_of_pair_of_groups} for a later account of marginal-series theory for pairs of groups.

A later group-theoretic treatment specifically devoted to nilpotency of pairs of groups is given in \cite{hassanzadeh_pourmirzaei_kayvanfar_2013_on_the_nilpotency_of_a_pair_of_groups}.
\end{HisNote}


\subsection{Relative central series and nilpotent homomorphisms}
\label{subsection in appendix: relative nilpotency for groups}


\begin{dfn}
\label[dfn]{definition: relative central series}
Let $G$ be a group and $N\trianglelefteq G$ be a normal subgroup. A \emph{$G$-relative central series} for $N$ is a finite sequence of normal subgroups of $G$,
\[ 1=N_{0}\subseteq N_{1}\subseteq\cdots\subseteq N_{c}=N, \]
such that $[G, N_{i+1}]\subseteq N_{i}$ for every $0\leq i<c$, or equivalently, $ N_{i+1}/ N_{i}\subseteq Z(G/ N_{i})$.

The normal subgroup $N$ is called \emph{nilpotent relative to $G$}, or \emph{$G$-nilpotent}, if it admits such a series. We define its \emph{$G$-nilpotency class}, denoted by $\ncl_G(N)$, as the least possible length. For notational convenience, we set $\ncl_G(N)= \infty$ when \(N\) is not \(G\)-nilpotent. When $N=G$, we simply say that $G$ is \emph{nilpotent} and write $\ncl(G)\coloneqq\ncl_G(G)$.
\end{dfn}


\begin{dfn}
\label[dfn]{definition: relative lower and upper central series}
Let $G$ be a group and $N\trianglelefteq G$ be a normal subgroup.
\begin{enumerate}
    \item The \emph{$G$-relative lower central series} of $N$ is defined inductively by $\relLCS_{0}(G,N)\coloneqq N$ and $\relLCS_{i+1}(G,N)\coloneqq[G,\relLCS_{i}(G,N)]$:
    \[ \cdots \subseteq \relLCS_{i}(G,N) \subseteq \cdots\subseteq \relLCS_{1}(G,N)=[G,N] \subseteq \relLCS_{0}(G,N) =N. \]
    \item The \emph{$G$-relative upper central series} of $N$ is defined inductively by $Z_0(G, N)\coloneqq1$ and $Z_{i+1}(G, N)\coloneqq \{x\in N\mid[G,x]\subseteq Z_i(G, N)\}$:
    \[ Z_{0}(G, N)=1 \subseteq Z_{1}(G, N) =N\cap Z(G) \subseteq \cdots\subseteq Z_{i}(G, N) \subseteq \cdots. \]
\end{enumerate}
Note that we use an indexing convention that starts in degree $0$ in the definition of the lower central series.
When $N=G$, we write $\grpLCS_{i}(G)\coloneqq\relLCS_{i}(G,G)$ and $Z_i(G)\coloneqq Z_{i}(G,G)$; these are the usual lower and upper central series of $G$.
\end{dfn}

In fact, if the $G$-relative lower and upper central series have finite length, then they are $G$-relative central series.

\begin{lem}
\label[lem]{proposition: relative upper central series as intersection}
For every $i\geq0$, we have $Z_i(G, N)=N\cap Z_i(G)$.
\end{lem}

\begin{proof}
For $x\in N$, we have $[G,x]\subseteq N$ by normality. Therefore, if the equality holds in degree $i$, $x\in Z_{i+1}(G, N)$ if and only if $[G, x]\subseteq Z_i(G)$. This completes the proof.
\end{proof}

\begin{thm}
\label[thm]{theorem: characterizations of nilpotent groups}
Let $G$ be a group and $N\trianglelefteq G$ be a normal subgroup, and let $c\geq0$. The following conditions are equivalent.
\begin{enumerate}[label={(\roman*)}]
    \item\label{item:G-nilpotent_of_class_c} The subgroup $N$ is $G$-nilpotent of class at most $c$.
    \item\label{item:relLCS_stops_at_c} $\relLCS_{c}(G,N)=1$.
    \item\label{item:relUCS_stops_at_c} $Z_{c}(G, N)=N$.
    \item\label{item:N_in_contained_in_absUCS_at_c} $N\subseteq Z_c(G)$.
\end{enumerate}
In particular, $\ncl_G(N)=\min\{i\geq 0\mid \Gamma_i(G,N)=1\}=\min\{i\geq 0\mid Z_i(G,N)=N\}$.
\end{thm}

\begin{proof}
\ref{item:relLCS_stops_at_c} or \ref{item:relUCS_stops_at_c} $\Rightarrow$ \ref{item:G-nilpotent_of_class_c}: immediate.
\ref{item:G-nilpotent_of_class_c} $\Rightarrow$ \ref{item:relLCS_stops_at_c}, \ref{item:relUCS_stops_at_c}: let $1=N_0\subseteq\cdots\subseteq N_c=N$ be any $G$-relative central series for $N$. Then, observe that induction gives $\relLCS_i(G,N)\subseteq N_{c-i}$ and $N_i\subseteq Z_i(G,N)$ for $0\leq i\leq c$. Thus we have $\relLCS_c(G,N)=1$ and $Z_c(G,N)=N$.
\ref{item:relUCS_stops_at_c} $\Leftrightarrow$ \ref{item:N_in_contained_in_absUCS_at_c}: follows from \cref{proposition: relative upper central series as intersection}.
\end{proof}

Taking $N=G$ in the preceding theorem gives the usual characterizations of nilpotent groups.

\begin{rem}
\label[rem]{remark: lower and upper operator central lengths}
For a general group with operators, the lower and upper operator-central series need not determine the same length, even when the relevant nilpotency conditions hold; see \cite[Example 4]{mohamed_1963_on_series_of_subgroups_related_to_groups_of_automorphisms}. The coincidence in \cref{theorem: characterizations of nilpotent groups} is a special feature of the conjugation action of an ambient group on a normal subgroup, as reflected in the identity $Z_i(G,N)=N\cap Z_i(G)$.
\end{rem}

\begin{eg}
    A normal subgroup $N$ of a group $G$ is nilpotent with $\ncl_G(N)=0$ if and only if $N$ is trivial; it is nilpotent with $\ncl_G(N)\leq1$ if and only if $N$ is centralized by $G$, which means $[G,N]=1$.
\end{eg}

\begin{dfn}
\label[dfn]{definition: nilpotent group homomorphism}
A surjective group homomorphism $\phi\colon G\twoheadrightarrow H$ is called \emph{nilpotent} if the kernel $\Ker(\phi)$ is $G$-nilpotent.
Its nilpotency class is denoted by $\ncl(\phi)\coloneqq \ncl_G(\Ker(\phi))$.
\end{dfn}

The characterization by iterated central extensions is now a direct relative analogue of the usual central-series description. Recall that a surjective group homomorphism $\pi\colon G\twoheadrightarrow H$ is a \emph{central extension} if $\Ker(\pi)\subseteq Z(G)$.

\begin{lem}
\label[lem]{lemma:central_series_correspond_to_central_extensions}
    Let $N\subseteq M$ be normal subgroups of $G$, and let $\pi\colon G/N \twoheadrightarrow G/M$ be the induced surjection. Then, the inclusion $N\subseteq M$ is $G$-central (i.e.\ $[G,M]\subseteq N$) if and only if $\pi$ is a central extension.
\end{lem}

\begin{proof}
    The condition $[G,M]\subseteq N$ is equivalent to $M/N\subseteq Z(G/N)$. Since $M/N=\Ker(\pi)$, this is equivalent to saying that $\pi$ is a central extension.
\end{proof}

\begin{thm}
\label[thm]{theorem: characterization of relative nilpotency}
Let $\phi\colon G\twoheadrightarrow H$ be a surjective group homomorphism and let $c\geq 0$. Then, $\phi$ is nilpotent of class at most $c$ if and only if $\phi$ admits a factorization as a composite of $c$ central extensions.
\end{thm}

\begin{proof}
($\Rightarrow$): Put $N\coloneqq\Ker(\phi)$. Let $1= N_{0}\subseteq\cdots\subseteq N_{c}=N$ be a $G$-relative central
series for $N$ (we may assume the length is $c$, by repeating terms if necessary). Then this sequence induces a factorization of $\phi$
\[\begin{tikzcd}
G=G/N_{0} \arrow[two heads]{r} &
G/N_{1} \arrow[two heads]{r} &
\cdots \arrow[two heads]{r} &
G/N_{c}\cong H
\end{tikzcd}\]
into $c$ central extensions by \cref{lemma:central_series_correspond_to_central_extensions}.
($\Leftarrow$): Conversely, from a factorization of $\phi$ into $c$ central extensions, take the kernels in $G$ of the successive composites. These kernels form a $G$-relative central series of length $c$, by \cref{lemma:central_series_correspond_to_central_extensions} again.
\end{proof}

From this theorem, we can also define a nilpotent homomorphism as a finite sequence of central extensions.
We will freely use this perspective in what follows, without further mention.

In particular, the lower and upper relative central series give two canonical factorizations of a nilpotent homomorphism. Let $\phi\colon G\twoheadrightarrow H$ be nilpotent of class at most $c$ and put $N\coloneqq\Ker(\phi)$. Then both sequences
\[\begin{tikzcd}
    G\cong G/\relLCS_{c}(G,N) \arrow[two heads]{r} &
    G/\relLCS_{c-1}(G,N) \arrow[two heads]{r} &
    \cdots \arrow[two heads]{r} &
    G/\relLCS_{0}(G,N)= H
\end{tikzcd}\]
and
\[\begin{tikzcd}
    G=G/Z_0(G, N) \arrow[two heads]{r} &
    G/ Z_{1}(G, N) \arrow[two heads]{r} &
    \cdots \arrow[two heads]{r} &
    G/Z_{c}(G, N)\cong H
\end{tikzcd}\]
are factorizations into central extensions.

\begin{eg}
    A surjective group homomorphism $\phi\colon G \twoheadrightarrow H$ is nilpotent with $\ncl(\phi)=0$ if and only if $\Ker(\phi)$ is trivial, and hence $\phi$ is an isomorphism; it is nilpotent with $\ncl(\phi)\leq1$ if and only if $\Ker(\phi)$ is contained in $Z(G)$, i.e., $\phi$ is a central extension.
\end{eg}

%For the terminal homomorphism $G\twoheadrightarrow1$, the preceding theorem says that $G$ is nilpotent of class at most $c$ if and only if $G\twoheadrightarrow1$ is a composite of $c$ central extensions. Thus all of the ordinary statements from the absolute theory are contained in the relative formulation.

The relative lower central series lies between the ordinary lower central series of the normal subgroup and that of the ambient group.

\begin{prop}
\label[prop]{proposition: comparison of lower central series}
    For a normal subgroup $N\subseteq G$, the inclusions $\Gamma_i(N) \subseteq \Gamma_i(G, N)\subseteq \Gamma_i(G)$ hold for all $i \geq 0$.
\end{prop}

\begin{proof}
When $i=0$, we have $\Gamma_0(N)=\Gamma_0(G,N)=N \subseteq G=\Gamma_0(G)$. If the inclusions hold for $i\geq 0$, then it follows that $\Gamma_{i+1}(N) = [N,\Gamma_i(N)] \subseteq [G,\Gamma_i(N)] \subseteq [G,\Gamma_i(G,N)]=\Gamma_{i+1}(G,N)$ and $\Gamma_{i+1}(G,N)=[G,\Gamma_i(G,N)] \subseteq [G,\Gamma_i(G)]=\Gamma_{i+1}(G)$.
\end{proof}

\begin{cor}
\label[cor]{corollary: relative nilpotency implies ordinary nilpotency}
    \begin{enumerate}
        \item If a normal subgroup $N\subseteq G$ is $G$-nilpotent of class at most $c$, then $N$ is nilpotent of class at most $c$.
        \item If $G$ is nilpotent of class at most $c$, then every normal subgroup $N$ of $G$ is $G$-nilpotent of class at most $c$.
    \end{enumerate}
\end{cor}

\begin{rem}
\label[rem]{remark: ordinary nilpotency does not imply relative nilpotency}
    Absolute nilpotency does not imply relative nilpotency.
    For example, consider the symmetric group $S_3$ and the alternating group $A_3\trianglelefteq S_3$ of degree $3$. Then $A_3$ is abelian, hence nilpotent. But $[S_3,A_3]=A_3$, so $\relLCS_{i}(S_3,A_3)=A_3$ for every $i\geq0$.
\end{rem}

\begin{lem}
\label[lem]{lemma: lower central series with cyclic quotient}
Let $N\trianglelefteq G$ be a normal subgroup such that $G/N$ is cyclic. Then, we have $\grpLCS_i(G)=\relLCS_i(G,N)$ for every $i\geq 1$. In particular, $N$ is $G$-nilpotent if and only if $G$ is nilpotent.
\end{lem}

\begin{proof}
First, we prove the case with $i=1$. The inclusion $[G,N] \subseteq [G,G]$ is trivial. The subgroup $N/[G,N]$ is central in $G/[G,N]$, and the corresponding quotient is isomorphic to $G/N$, which is cyclic. Then we claim that $G/[G,N]$ is abelian; indeed, putting $K=N/[G,N]$ and $H=G/[G,N]$, there is some $h \in H$ such that $H/K=\langle \overline{h}\rangle$. For $x,y \in H$, we can write $x=h^m k_1$ and $y=h^nk_2$, where $k_1,k_2\in K$ and $m,n\in \ZZ$. Since $k_1,k_2$ are in the center of $H$, we have $xy=h^mk_1h^n k_2 = h^{m+n}k_1k_2 = h^nk_2 h^mk_1 = yx$, and hence $H=G/[G,N]$ is abelian.
From this, it follows that $[G,G]\subseteq[G,N]$, and therefore $\grpLCS_1(G)=\relLCS_1(G,N)$. The equality in every degree
$i\geq1$ follows by induction.
%If $H$ is $G$-nilpotent of positive class, the equality of the two series shows that $G$ is nilpotent. If $H=1$, then $G$ is cyclic. The converse follows from the same equality.
\end{proof}

We record the following proposition, which will be used repeatedly in the main text.

\begin{prop}
\label[prop]{proposition: functoriality of relative lower central series}
Let $\phi\colon G\twoheadrightarrow H$ be a surjective homomorphism of groups. Let $N\trianglelefteq G$ be a normal subgroup, and put $M\coloneqq\phi(N)$, which is also normal by the surjectivity of $\phi$. Then, $\phi(\relLCS_i(G, N)) = \relLCS_i(H, M)$ for every $i\ge0$.
%In particular, $\ncl_H(M)\le\ncl_G(N)$.
\end{prop}

\begin{proof}
    We use the fact that $\phi([A,B])=[\phi(A),\phi(B)]$ in general. For $i=0$, we have $\phi(\relLCS_0(G,N))=\phi(N)=M=\relLCS_0(H,M)$. Suppose $\phi(\relLCS_i(G, N)) = \relLCS_i(H, M)$ for $i\geq 0$. Then, $\phi(\relLCS_{i+1}(G, N)) = \phi([G,\relLCS_i(G,N)]) =[H,\phi(\relLCS_i(G,N))]= [H,\relLCS_i(H,M)]=\relLCS_{i+1}(H, M)$. Thus, induction completes the proof.
\end{proof}

The lower series also gives the universal approximation by homomorphisms of bounded nilpotency class.

\begin{thm}
\label[thm]{theorem: universal relative nilpotent truncation}
Let $\phi\colon G\twoheadrightarrow H$ be a surjective group homomorphism and let $c\ge0$. With the abbreviation $\relLCS_c= \relLCS_c(G,\Ker(\phi))$, the induced homomorphism $\phi_c\colon G/\relLCS_c \twoheadrightarrow H$ is nilpotent of class at most $c$.

Moreover, the factorization $\phi=\phi_c\circ \pi$, where $\pi\colon G \to G/\relLCS_c$ is the quotient map, is universal among factorizations of $\phi$ whose second homomorphism has class at most $c$. That is, if $\phi=\theta\circ \xi$ with $\xi\colon G\twoheadrightarrow K$ surjective and $\theta\colon K\twoheadrightarrow H$ nilpotent of class at most $c$, then $\xi$ factors uniquely through $G/\relLCS_{c}$.
\end{thm}

\begin{proof}
One can easily check that $\pi(\Ker(\phi))=\Ker(\phi_c)$. By \cref{proposition: functoriality of relative lower central series}, we have $\relLCS_c(G/\relLCS_c,\Ker(\phi_c)) = \pi(\relLCS_c(G,\Ker(\phi)))=1$, which shows that $\phi_c$ is nilpotent of class at most $c$.
For the universal property, $\xi(\relLCS_{c}(G,\Ker(\phi))) =\relLCS_{c}(K,\Ker(\theta))=1$, so $\xi$ factors uniquely through the quotient.
\end{proof}


\begin{dfn}
\label[dfn]{definition: relative nilpotent truncation}
Let $c\geq 0$. For a surjective group homomorphism $\phi\colon G\twoheadrightarrow H$, we refer to the induced nilpotent homomorphism $\phi_c\colon G/\relLCS_{c}\twoheadrightarrow H$ as the \emph{$c$-nilpotent truncation} of $\phi$.
\end{dfn}


\subsection{Stable properties}
\label{subsection:stable_property_of_nilpotency}


\begin{prop}
\label[prop]{proposition:nilpotency_for_group_descend_along_surjection}
    Consider a commutative diagram of surjective homomorphisms
    \[\begin{tikzcd}
        G' \arrow[two heads]{r}{\theta} \arrow[two heads]{d}[swap]{\phi'} & G \arrow[two heads]{d}{\phi} \\
        H' \arrow[two heads]{r}[swap]{\kappa} & H\rlap{.}
    \end{tikzcd}\]
    with $\theta(\Ker(\phi'))=\Ker(\phi)$. If $\phi'$ is nilpotent of class at most $c$, then so is $\phi$.
\end{prop}

\begin{proof}
    It follows from \cref{proposition: functoriality of relative lower central series} that $\relLCS_c(G,\Ker(\phi)) = \theta(\relLCS_c(G',\Ker(\phi')))=1$, provided that $\relLCS_c(G',\Ker(\phi'))=1$.
\end{proof}


\begin{prop}
\label[prop]{proposition: relative lower central series under pullback}
Consider a pullback square of a surjective homomorphism $\phi$:
\[\begin{tikzcd}
    G' \arrow{r}{\theta} \arrow[two heads]{d}[swap]{\phi'} & G \arrow[two heads]{d}{\phi} \\
    H' \arrow{r}[swap]{\kappa} & H\arrow[phantom]{lu}[very near end]{\lrcorner} \rlap{.}
\end{tikzcd}\]
If $\phi$ is nilpotent of class at most $c$, then so is $\phi'$.
Moreover, if $\kappa$ is surjective, then the converse also holds; in particular, $\ncl(\phi)=\ncl(\phi')$ holds.
\end{prop}



\begin{proof}
Note that $\theta$ induces an isomorphism $\Ker(\phi')\cong \Ker(\phi)$ between the kernels.
By \cref{proposition: functoriality of relative lower central series}, it also induces $\relLCS_i(G',\Ker(\phi'))\cong \theta(\relLCS_i(G',\Ker(\phi'))) = \relLCS_i(\Image(\theta),\Ker(\phi))$ for every $i\geq 0$.
Since $\relLCS_i(\Image(\theta),\Ker(\phi)) \subseteq \relLCS_i(G,\Ker(\phi))$, if $\phi$ is nilpotent of class $c$, then so is $\phi'$.

The last assertion follows from \cref{proposition:nilpotency_for_group_descend_along_surjection}.
\end{proof}


\begin{lem}
\label[lem]{lemma: relative lower central series under subgroups}
Let $\phi\colon G\twoheadrightarrow H$ and $\phi'\colon G'\twoheadrightarrow H'$ be two surjective group homomorphisms. Suppose that we have inclusions $\iota_G\colon G'\hookrightarrow G$ and $\iota_H\colon H'\hookrightarrow H$ such that $\phi\circ \iota_G = \iota_H\circ \phi'$. If $\phi$ is nilpotent of class at most $c$, then so is $\phi'$.
\end{lem}

\begin{proof}
Put $N=\Ker(\phi)$ and $N'=\Ker(\phi')$. Then we have an inclusion $N' \subseteq N$. Regarding $G'\subseteq G$ as a subgroup via $\iota_G$, we obtain $\relLCS_{i}(G',N') \subseteq \relLCS_{i}(G,N)$ for every $i \geq 0$ inductively. Therefore, the assertion follows from it.
\end{proof}


\begin{prop}
\label[prop]{proposition: relative nilpotency under finite products}
Let $\phi\colon G\twoheadrightarrow H$ and $\phi'\colon G'\twoheadrightarrow H'$ be two surjective group homomorphisms. Then, the product $\phi\times \phi'$ is nilpotent of class at most $c$ if and only if $\phi$ and $\phi'$ are nilpotent of class at most $c$.
In particular, $\ncl(\phi\times \phi') = \max\{\ncl(\phi),\ncl(\phi')\}$ holds.
\end{prop}

\begin{proof}
We use the fact that $[A\times A',B\times B'] = [A,B]\times [A',B']$ as subsets of $G\times G'$ for subgroups $A,B\subseteq G$ and $A',B'\subseteq G'$.
Put $N=\Ker(\phi)$ and $N'=\Ker(\phi')$.
Then, by induction, we obtain $\relLCS_i(G\times G',N\times N') \cong \relLCS_i(G,N)\times \relLCS_i(G',N')$ for every $i\geq 0$.
Since $\Ker(\phi\times\phi')=N\times N'$, the assertion easily follows from it.
\end{proof}





\begin{prop}
\label[prop]{proposition: composition of nilpotent group homomorphisms}
Let $\phi\colon G\twoheadrightarrow H$ and $\psi\colon H \twoheadrightarrow K$ be surjective group homomorphisms.
\begin{enumerate}
    \item\label{item:composition_of_nilpotent_hom} If $\phi$ is nilpotent of class at most $c$ and $\psi$ is nilpotent of class at most $d$, then $\psi\circ \phi$ is nilpotent of class at most $c+d$.
    \item\label{item:cancelation_of_nilpotent_hom} If $\psi\circ\phi$ is nilpotent of class at most $c$, then both $\phi$ and $\psi$ are so.
\end{enumerate}
Consequently, whenever all three homomorphisms are nilpotent,
\[ \max\{\ncl(\phi),\ncl(\psi)\} \leq \ncl(\psi\circ\phi) \leq \ncl(\phi)+\ncl(\psi).\]
\end{prop}

\begin{proof}
\eqref{item:composition_of_nilpotent_hom}: When $\phi$ (resp.\ $\psi$) is a composite of $c$ (resp.\ $d$) central extensions, $\psi\circ\phi$ is a composite of $(c+d)$ central extensions. Hence, it follows from \cref{theorem: characterization of relative nilpotency}.

\eqref{item:cancelation_of_nilpotent_hom}: Put $\chi=\psi\circ\phi$ and assume that $\chi$ is nilpotent of class $c$. Then, $\Ker(\phi) \subseteq \Ker(\chi)$ gives $\relLCS_c(G,\Ker(\phi)) \subseteq \relLCS_c(G,\Ker(\chi))=1$, which shows that $\phi$ is nilpotent of class at most $c$.
Also, because $\phi(\Ker(\chi))=\Ker(\psi)$, \cref{proposition:nilpotency_for_group_descend_along_surjection} implies that $\psi$ is nilpotent of class at most $c$.
\end{proof}




\begin{comment}
\begin{prop}
\label[prop]{proposition: relative nilpotency in extensions}
\label[prop]{proposition: composition of nilpotent group homomorphisms}
    Let $K\subseteq N$ be normal subgroups of a group $G$ and let $f\colon G\twoheadrightarrow H$ and $g\colon H\twoheadrightarrow L$ be surjective group homomorphisms. 
    Then,
    \begin{align*}
        \max\{\ncl_G(K),\ncl_{G/K}(N/K)\}
        &\leq \ncl_G(N)
        \leq \ncl_G(K)+\ncl_{G/K}(N/K),\\
        \max\{\ncl(f),\ncl(g)\}
        &\leq \ncl(g\circ f)
        \leq \ncl(f)+\ncl(g).
    \end{align*}
    In particular, the following hold.
    \begin{enumerate}
        \item\label{item: relative nilpotency is closed under composition} If $f$ and $g$ are nilpotent, then so is $g\circ f$.
        \item\label{item: relative nilpotency descends to factors} If $g\circ f$ is nilpotent, then both $f$ and $g$ are nilpotent.
    \end{enumerate}
\end{prop}

\begin{proof}
    The inclusions $\Gamma_i(G,K)\subseteq\Gamma_i(G,N)$ and the functoriality under $G\twoheadrightarrow G/K$ give the first inequality. For the second, put $c\coloneqq\ncl_G(K)$ and $d\coloneqq\ncl_{G/K}(N/K)$, and suppose that $c,d<\infty$. Then, $\Gamma_d(G,N)\subseteq K$, and hence $\Gamma_{c+d}(G,N)=\Gamma_c(G,\Gamma_d(G,N))\subseteq\Gamma_c(G,K)=1$.

    For the homomorphism statement, put $M\coloneqq\Ker(f)$ and $N'\coloneqq\Ker(g\circ f)$. Under $G/M\cong H$, the subgroup $N'/M$ identifies with $\Ker(g)$. Applying the first assertion to $M\subseteq N'$ gives the desired inequalities, from which \eqref{item: relative nilpotency is closed under composition} and \eqref{item: relative nilpotency descends to factors} follow.
\end{proof}

\begin{prop}
\label[prop]{proposition: relative nilpotency in extensions}
Let $K\subseteq N$ be normal subgroups of $G$. If $K$ is
$G$-nilpotent of class at most $c$ and $N/K$ is
$(G/K)$-nilpotent of class at most $d$, then $N$ is $G$-nilpotent of
class at most $c+d$.
\end{prop}

\begin{proof}
For the quotient map $\pi\colon G\twoheadrightarrow G/K$,
functoriality gives
$\pi(\relLCS_{d}(G,N))=\relLCS_{d}(G/K,N/K)=1$. Hence
$\relLCS_{d}(G,N)\subseteq K$, and taking a further $c$ relative
commutators gives $\relLCS_{c+d}(G, N)=1$.
\end{proof}

\end{comment}

We collect the closure properties for relatively nilpotent normal subgroups.

\begin{prop}
\label[prop]{proposition: products and intersections of relatively nilpotent subgroups}
    Let $N,M\trianglelefteq G$ be normal subgroups, and let $c, d\geq 0$. Then the following hold:
    \begin{enumerate}
        \item\label{item:subgroup_of_relative_nilpotent} If $N$ is $G$-nilpotent of class at most $c$ and $M\subseteq N$, then $M$ is $G$-nilpotent of class at most $c$.
        \item\label{item:intersection_of_relative_nilpotent} If $N$ and $M$ are $G$-nilpotent of class at most $c$ and $d$ respectively, then $N \cap M$ is $G$-nilpotent of class at most $\min\{c,d\}$. Hence, $\ncl_G(N\cap M)\le\min\{\ncl_G(N), \ncl_G(M)\}$.
        \item\label{item:product_subgroup_of_relative_nilpotent} If $N$ and $M$ are $G$-nilpotent of class at most $c$ and $d$ respectively, then $NM$ is $G$-nilpotent of class at most $\max\{c,d\}$. In particular, we have $\ncl_G(NM) = \max\{\ncl_G(N), \ncl_G(M)\}$.
    \end{enumerate}
\end{prop}

\begin{proof}
\eqref{item:subgroup_of_relative_nilpotent}: By \cref{theorem: characterizations of nilpotent groups}, $N$ is nilpotent of class at most $c$ if and only if $N\subseteq Z_c(G)$. Therefore, $M$ is nilpotent of class at most $c$ when $M \subseteq N$.
\eqref{item:intersection_of_relative_nilpotent}: It follows immediately from \eqref{item:subgroup_of_relative_nilpotent}.
\eqref{item:product_subgroup_of_relative_nilpotent}: Put $r=\max\{c,d\}$. Then $N,M \subseteq Z_r(G)$, which implies that $NM \subseteq Z_r(G)$. Hence $NM$ is $G$-nilpotent of class at most $r$. In particular, $\ncl_G(NM) \leq \max\{\ncl_G(N),\ncl_G(M)\}$. The reverse inequality follows from \eqref{item:subgroup_of_relative_nilpotent}, and hence equality holds.
\end{proof}

